\thispagestyle{plain}
\begin{center} \begin{huge} 摘要 \end{huge} \end{center}
\noindent
過去，工業控制系統（Industrial Control Systems, ICS）多被視為封閉且與外部網路隔離的環境。然而，隨著資訊科技（Information Technology, IT）與營運科技（Operational Technology, OT）網路逐漸融合，ICS 面臨的外部威脅也隨之增加。由於 Modbus TCP 等傳統工業通訊協定缺乏原生安全機制，難以滿足現代資安防護需求，使關鍵基礎設施暴露於更高的攻擊風險。攻擊者可能利用協定弱點注入未授權指令、干擾物理流程，或繞過邊界防禦，進而影響系統穩定性與營運安全。同時，工業現場對即時性與穩定性具有高度要求，因此如何在不增加通訊延遲與系統負擔的前提下導入有效監測機制，成為 ICS 資安防護的重要課題。

為解決上述問題，本論文提出一個整合 ISA/IEC 62443-3-3 Foundational Requirements（FRs）與 ICS Cyber Kill Chain 的 Modbus TCP 威脅監測架構，用以識別工業控制環境中的異常通訊與攻擊行為。本研究採用 Suricata 與 Graylog 建構開源式被動監測系統，透過網路流量分析、偵測規則與日誌關聯機制，驗證不同攻擊情境下的威脅監測能力。此架構以開源工具為基礎，可降低部署與維護成本，使資安預算有限的企業亦能提升工業網路的安全可視性。

此外，當系統偵測到可疑事件時，會透過 Slack 即時通知維運人員，以支援快速事件回應。本研究亦將告警事件對應至 MITRE ATT\&CK for ICS 技術項目，協助維運人員理解攻擊行為所屬的戰術階段與技術特徵。

本研究在虛擬化實驗室平台中，透過概念驗證（PoC）實驗驗證了該框架的防禦能力，測試場景涵蓋偵察掃描、協定模糊測試、未授權參數竄改、關鍵診斷指令注入，以及拒絕服務（DoS）洪水攻擊。實證結果顯示，本系統具備優異的通報即時性，所有攻擊情境的平均端到端告警延遲均保持在 10 秒以下。在峰值壓力測試的高吞吐量負載下，該被動防禦管線成功承受了每秒高達 29,128.00 個封包（PPS，約 13.36 Mbps）的持續洪水轟炸，並實現了 Suricata 核心 Socket 層 0.00\% 丟包、以及 Graylog SIEM 日誌收容層 0.00\% 漏日誌的實測表現。此外，透過 CICModbusDataset2023 與 4SICS Geek Lounge 兩個第三方公開工控資料集的交叉驗證，進一步證實了本研究自訂規則在異質外部 OT 網路環境中的可移植性與偵測效能。
\newpage
整體而言，本研究所提出所提出的框架在不主動干擾即時 Modbus TCP 通訊流程前提下，提升了工業控制網路中的安全可視性、告警即時性與事件追溯能力，證實其部署於輕量化、消費級硬體上的高度實務可行性。

{關鍵字: Modbus TCP、ISA/IEC 62443-3-3、工控阻殺鏈、Suricata、Graylog、MITRE ATT\&CK for ICS}
\newpage